\chapter{Аудит топологических квантов четырёхмерного объёма}
\label{ch:v10-cell-birth-four-volume-topological-quantum-audit}

\section{Проверяемые кандидаты}

Спектральный счёт оставил свободным элементарный объём. Проверим восемь
целочисленных кандидатов: характеристику Эйлера, сигнатуру, число
Понтрягина, индекс Дирака, число обвитий, число симплексов, индекс потока
GNVW и размерность $K_{43}$.

Каждый кандидат оценивается по пяти условиям: целочисленная защищённость,
внутреннее происхождение, размерность $L^4$, отображение в физический объём
и разрыв масштабной орбиты. Матрица имеет ранг $2$, а вектор полных
проходов равен
\begin{equation}
 (0,0,0,0,0,0,0,0)^T.
 \label{eq:v10-topological-volume-pass-vector}
\end{equation}
Все кандидаты защищены и внутренни, но ни один не имеет размерности $L^4$
и ни один не разрывает физическую масштабную орбиту.

\section{Кратность не равна элементарному объёму}

Топология может дать целое $n$, после чего полный объём имеет вид
\begin{equation}
 V=n v_0.
 \label{eq:v10-topological-volume-factorization}
\end{equation}
Однако при $v_0\mapsto s^4v_0$ целое $n$ не меняется, а объём масштабируется.
Топологическая плотность
\begin{equation}
 \rho_{\rm top}=\frac{n}{V},
 \qquad \rho_{\rm top}V=n
 \label{eq:v10-topological-density}
\end{equation}
получает размерность от метрики и преобразуется как $s^{-4}$; она не
создаёт метрику сама.

Родитель, фиксирующий целую кратность, имеет в координатах отклонения
кратности и логарифма $v_0$ гессиан
\begin{equation}
 \begin{pmatrix}1&0\\0&0\end{pmatrix},
 \qquad \rank=1,
 \qquad \dim\ker=1.
 \label{eq:v10-topological-volume-parent-hessian}
\end{equation}
Его ядро совпадает с направлением элементарного объёма.

\begin{theorem}[Топология выбирает кратность, но не физический квант]
Все проверенные топологические инварианты могут ограничивать целые числа и
классы носителей, однако физический объём неизбежно факторизуется как
$V=n v_0$ с невыведенным размерным множителем $v_0$.
\end{theorem}

\begin{nogocriterion}[Запрет отождествления индекса с объёмом]
Целое топологическое число нельзя объявлять величиной размерности $L^4$.
Для такого перехода необходим отдельно выведенный коэффициент объёмной
плотности или размерная мера.
\end{nogocriterion}

\begin{protocol}[Статус топологических кандидатов]\begin{enumerate}\raggedright
\item проверенных кандидатов: \textbf{$8$};
\item полных проходов: \textbf{$0/8$};
\item ранг матрицы кандидатов: \textbf{$2$};
\item топологическая кратность: \textbf{выводима};
\item элементарный четырёхмерный объём: \textbf{$0/1$};
\item абсолютная энергия часов: \textbf{$0/1$};
\item следующий гейт: \textbf{родитель объёмной плотности кривизны}.
\end{enumerate}\end{protocol}

Машинный сертификат сохранён в
\path{s2t_v10_cell_birth_four_volume_topological_quantum_candidate_audit_gate_results.json}.